% Recuperación de 02_trit_geometrias y de la holonomía con memoria.
\subsubsection{Temporal parametrization of compatibility and tritic orientation}
\label{trit-temporal-compatibilidad}

The temporal interpretation of the TRIT is established on transitions of the
complete state. Let \(x\xrightarrow{\rm adm}y\) be the admissibility relation determined
by the TPK operations on \(\widehat X\). A trajectory parametrized
by a discretely ordered set \(\mathcal T\) is a map
\[
 \gamma:\mathcal T\longrightarrow\widehat X,\qquad
 \gamma(t)\xrightarrow{\rm adm}\gamma(t^+),
\]
where \(t^+\) is the successor of \(t\). The compatibility relation precedes
the choice of parameter: the order expresses its transitions as
an evolution. When the term section is used, the
projection \(p:E\to\mathcal T\) and condition
\(p\circ\gamma=\operatorname{id}_{\mathcal T}\) are also declared. This specifies
the state space, temporal order and information transported.

The signature \(\tau\) selects a quadratic regime and admits a
temporal readout relative to that parametrization. Type \(+1\), with
\(J_{+1}^{2}=-I\), represents return and memory; type \(0\), with
\(J_0^{2}=0\), represents the transition threshold; type \(-1\), with
\(J_{-1}^{2}=I\), represents opening in two eigendirections. In the
model's temporal terminology, these functions correspond,
respectively, to past, present and future. The correspondence concerns
the operational position of a transition within a history:
preserved antecedent, update and admissible continuation.

This interpretation acquires content through three distinct operations.
Restricting a trajectory recovers the prefix already constructed. The
update \(\mathcal U_t\) composes selection, transport and registration in
the observed transition. Prolongation considers the continuations that
preserve the boundary conditions of that prefix. If
\(\mathcal P_n(x)\) denotes the set of admissible continuations of
length \(n\) from \(x\), deletion of the last step defines
\[
 r_n:\mathcal P_{n+1}(x)\longrightarrow\mathcal P_n(x).
\]
An unbounded continuation is a collection \((\gamma_n)_{n\geq1}\) satisfying
\(r_n(\gamma_{n+1})=\gamma_n\). The inverse system thus preserves
the restrictions that each depth imposes on the preceding
ones. Existence of a compatible collection is established in the
corresponding prolongation domain; merely setting up the
sets \(\mathcal P_n(x)\) does not replace that proof.

Bidirectionality therefore combines reconstruction of antecedents
with compatibility of continuations. On the signature sequence, the
involution \(C_{\rm rev}\) reverses the order and changes the sign of each
component. The trajectory is also acted on by the transformations of
position, sheet and boundary specified by transport. The reversed
signature is an observation of the conjugate trajectory, not a replacement
for its remaining coordinates. The threshold remains distinguished: the
visible value \(0\) retains its transition function even when its quotient
or historical register is nonzero.

The relationship with exponential geometries is equally precise. In
each fixed regime, \(\exp(tJ_\tau)\) composes parameters and preserves the
algebraic norm. Elliptic closure belongs to this local representation.
The phase return of \(\Gamma_9\), by contrast, acts on the trajectory
with memory, according to \eqref{eq:holonomia-memoria}. A trajectory may
recover its phase while retaining a new antecedent: its historical data have
changed even though the phase observation is the same. The parabolic threshold
retains a linear displacement; hyperbolic opening separates two
eigendirections of expansion and contraction. These three realizations
provide composition laws, while the TPK determines their
selection and effective concatenation.

The aperiodic suspension of Section \ref{sec:generacion} will make
this distinction explicit through a finite phase factor, an irrational rotation
and a memory degree. The connection between TRIT and temporality is thus
expressed through operations on trajectories and their quadratic
representations. Its content goes beyond a classification of signs: it preserves the
regime, direction of transport and relationship between antecedent,
transition and continuation.
